\documentclass[12pt, a4paper]{article}

\usepackage[utf8]{inputenc}
\usepackage[T1]{fontenc}
\usepackage{lmodern}
\usepackage{amsmath, amssymb, amsthm}
\usepackage{booktabs}
\usepackage{graphicx}
\usepackage{hyperref}
\usepackage{natbib}
\usepackage{geometry}
\usepackage{setspace}
\usepackage{caption}
\usepackage{array}
\usepackage{bm}

\geometry{margin=1in}
\onehalfspacing

\title{GSPQR: A Geometric Soft-Pullback Mechanism for\\
Quantile Interval Prediction in A-Share Markets}

\author{Liu Bowei\\
\small Hunan University of Information Technology\\
\small \texttt{20070316lbw@gmail.com}}

\date{April 2026}

\begin{document}

\maketitle

\begin{abstract}
Accurate uncertainty quantification is essential for risk-aware decision-making in
equity markets, yet existing interval prediction methods often treat interval
construction and point forecasting as disjoint tasks. Conformal prediction methods
such as EnbPI provide distribution-free coverage guarantees but produce conservative
intervals that do not adapt to the local confidence of the underlying predictor.
Standard quantile regression anchors intervals on the conditional median, discarding
the complementary directional signal from mean-based estimators. We propose GSPQR
(Geometric Soft-Pullback Quantile Regression), a lightweight mechanism that bridges
these approaches via an uncertainty-aware exponential gating function, formalizing
interval re-centering as an adaptive problem that interpolates between the conditional
mean and median based on the local signal-to-noise ratio. Evaluation on the HS300
constituent stocks (2023--2024, $n = 120{,}512$) at a nominal 80\% coverage level
shows that GSPQR achieves a Winkler Score of 0.4784, a 5.9\% improvement over EnbPI
($t = -68.5$, $p \ll 0.001$), while reducing interval width by 17.0\%. The $\beta$
parameter traces a convex Pareto frontier between calibration and cross-sectional
ranking ability, identifying an optimal balance at $\beta^* = 1.0$ under
regime-shift-intensive market conditions.
\end{abstract}

\textbf{Keywords:} quantile regression, interval prediction, conformal prediction,
LightGBM, A-share market, uncertainty quantification\\

\textbf{JEL Classification:} C53, C58, G17, G32

\hrule
\vspace{1em}

%%%%%%%%%%%%%%%%%%%%%%%%%%%%%%%%%%%%%%%%%%%%%%%%%%%
\section{Introduction}
%%%%%%%%%%%%%%%%%%%%%%%%%%%%%%%%%%%%%%%%%%%%%%%%%%%

Predicting future returns in equity markets is intrinsically a task of uncertainty
management. While point forecasts offer a singular view, interval prediction --- the
construction of a range $[\hat{l}_t, \hat{u}_t]$ containing the true return with a
specified probability --- provides the necessary granularity for position sizing, risk
management, and strategy evaluation. However, financial return distributions are
notoriously heavy-tailed, skewed, and non-stationary, rendering classical Gaussian
intervals ineffective. These challenges are especially acute in the Chinese A-share
market, which combines high retail participation, frequent policy-driven shocks, and
pronounced momentum-reversal dynamics.

The current state-of-the-art involves two primary paradigms. Conformal prediction
\citep{vovk2005} and its time-series adaptation EnbPI \citep{xu2021} offer
distribution-free coverage guarantees, but construct intervals by scaling a fixed
nonconformity score distribution without exploiting the predictive signal of the
underlying model. This leads to conservative, wide intervals: on our HS300 test set,
EnbPI produces a mean width of 0.3851, compared to 0.3197 for our method --- a 17.0\%
excess that directly reduces capital efficiency. Quantile regression
\citep{koenker1978}, implemented via gradient boosting \citep{ke2017}, yields
asymmetric, adaptive intervals but anchors them on the conditional median
$\hat{q}_{0.5}$, effectively ignoring the directional signal from mean-based predictors.
In the A-share market, where retail-driven sentiment and frequent regime shifts decouple
the mean and median, this choice is consequential: the pure median anchor achieves a
RankIC of only 0.0024 in our experiments.

\textbf{This paper introduces GSPQR (Geometric Soft-Pullback Quantile Regression)},
which treats interval construction as an adaptive center estimation problem. We propose
an exponential gating function that dynamically interpolates between the conditional
mean and median as a function of the local signal-to-noise ratio. Inspired by the
geometry of the logarithmic spiral, the mechanism preserves interval width while
optimizing placement. We make the following contributions:

\begin{enumerate}
    \item \textbf{Geometric soft-pullback.} A novel uncertainty gate based on the
    exponential decay of the mean-median deviation normalized by the confidence radius,
    with interpretable limiting behavior recovering the MSE baseline under high
    uncertainty and the Q50 anchor under high confidence.

    \item \textbf{Convex $\beta$ trade-off.} Empirical demonstration that the Winkler
    Score as a function of $\beta$ is convex, revealing a fundamental tension between
    interval calibration and cross-sectional ranking ability navigated by a single
    hyperparameter.

    \item \textbf{Strong empirical performance.} A Winkler Score of 0.4784 on the
    2023--2024 HS300 test set ($n = 120{,}512$), representing a 5.9\% improvement over
    EnbPI ($t = -68.5$, $p \ll 0.001$) and a 0.86\% MAE reduction relative to the
    MSE-center baseline despite no direct MSE optimization.
\end{enumerate}

\paragraph{Code Availability.}
The implementation of GSPQR and all experiment scripts are publicly available at:
\url{https://github.com/20070316lbw-netizen/snail-shell.git}.

\paragraph{Code Availability.}
The implementation of GSPQR, including data processing, model training, and evaluation pipelines, is publicly available at:

\url{https://github.com/20070316lbw-netizen/snail-shell.git}.

The remainder of the paper is organized as follows. Section~2 reviews related work.
Section~3 presents the GSPQR methodology. Section~4 reports experimental results.
Section~5 discusses findings and limitations. Section~6 concludes.

%%%%%%%%%%%%%%%%%%%%%%%%%%%%%%%%%%%%%%%%%%%%%%%%%%%
\section{Related Work}
%%%%%%%%%%%%%%%%%%%%%%%%%%%%%%%%%%%%%%%%%%%%%%%%%%%

\subsection{Quantile Regression and Gradient Boosting}

Quantile regression, introduced by \citet{koenker1978}, provides a robust framework
for estimating conditional quantiles without assuming a specific noise distribution,
by minimizing the asymmetric Pinball Loss. This robustness is particularly valuable
in financial return prediction, where heavy tails and outliers are prevalent. In modern
tabular settings, LightGBM \citep{ke2017} has become the de-facto standard for this
task, combining high-dimensional feature handling with native quantile objective support
and low computational overhead.

\subsection{Conformal Prediction for Time Series}

Conformal prediction \citep{vovk2005} provides distribution-free coverage guarantees
by calibrating interval width on a held-out set of nonconformity scores. For
non-stationary time series, EnbPI \citep{xu2021} adapts this framework using a
rolling calibration window. While EnbPI achieves valid marginal coverage, \citet{barber2023}
note that marginal coverage does not imply conditional coverage --- a distinction
especially consequential in cross-sectional equity settings where conditional
distributions vary substantially across stocks and regimes. Furthermore, conformal
methods calibrate width post-hoc without exploiting the predictive signal of the
underlying model, leading to systematically wider intervals than necessary.

\subsection{Interval Prediction in Financial Markets}

Accurate interval estimation is critical for portfolio risk management.
\citet{keilbar2022} demonstrate that standard parametric intervals frequently fail to
capture the systemic risk and heavy tails prevalent in financial returns, motivating
data-driven, asymmetric estimators that adapt to local distributional structure. Our
experiments corroborate this finding: the Residual baseline, which assumes stationary
Gaussian noise, achieves a coverage error of 0.1396 on the HS300 test set --- nearly
three times the tolerance of our composite scoring criterion.

\subsection{Positioning of GSPQR}

Existing methods treat interval construction and point prediction as separate problems.
GSPQR bridges this gap by introducing a geometric gate that continuously interpolates
between the two center estimates as a function of local uncertainty, yielding intervals
that are both well-calibrated and informationally efficient without modifying the
underlying predictors.

%%%%%%%%%%%%%%%%%%%%%%%%%%%%%%%%%%%%%%%%%%%%%%%%%%%
\section{Methodology}
%%%%%%%%%%%%%%%%%%%%%%%%%%%%%%%%%%%%%%%%%%%%%%%%%%%

\subsection{Problem Formulation}

Let $\mathcal{S} = \{s_1, \ldots, s_N\}$ denote the universe of $N$ stocks in the
HS300 index. At each trading day $t$, we observe a feature vector
$\mathbf{x}_{s,t} \in \mathbb{R}^{12}$ for stock $s$, and define the prediction target
as the one-month-ahead (20-trading-day) excess return $y_{s,t+h}$. The goal is to
construct a prediction interval $\hat{C}_{s,t} = [\hat{l}_{s,t},\, \hat{u}_{s,t}]$
satisfying:
%
\begin{equation}
\mathbb{P}\!\left(y_{s,t+h} \in \hat{C}_{s,t}\right) \geq 1 - \alpha
\end{equation}
%
where $\alpha = 0.2$ (nominal coverage 80\%). We evaluate interval quality using the
\textbf{Winkler Score} \citep{winkler1972}:
%
\begin{equation}
W_{s,t} = (\hat{u}_{s,t} - \hat{l}_{s,t})
+ \frac{2}{\alpha}\Bigl[
    \max(0,\, \hat{l}_{s,t} - y_{s,t+h})
  + \max(0,\, y_{s,t+h} - \hat{u}_{s,t})
\Bigr]
\end{equation}
%
and the \textbf{Coverage Error} $\text{CE} = |AC - (1-\alpha)|$, where $AC$ is the
empirical coverage. Lower values indicate better performance on both metrics.

\subsection{Base Quantile Predictor}

We train three independent LightGBM models at quantile levels
$\tau \in \{0.1, 0.5, 0.9\}$, each minimizing the Pinball Loss:
%
\begin{equation}
\mathcal{L}_\tau(\hat{y}, y) =
\begin{cases}
\tau \cdot (y - \hat{y}) & \text{if } y \geq \hat{y} \\
(1 - \tau) \cdot (\hat{y} - y) & \text{otherwise}
\end{cases}
\end{equation}
%
yielding predictions $\hat{q}_{0.1}$, $\hat{q}_{0.5}$, and $\hat{q}_{0.9}$ that form
the raw interval $[\hat{q}_{0.1},\, \hat{q}_{0.9}]$ with center $a_t = \hat{q}_{0.5,t}$
and confidence radius $r_t = \frac{1}{2}(\hat{q}_{0.9,t} - \hat{q}_{0.1,t})$.
A fourth LightGBM model trained under MSE produces the conditional mean estimate
$\hat{y}_t$. All four models share the same feature set, normalized via rolling
Z-score with a 60-day window.

All base predictors are trained independently for each configuration, 
yielding stable and consistent results across different $\beta$ settings.

\subsection{Interpretation as Adaptive Center Estimation}

The core hypothesis of GSPQR is that neither the mean nor the median alone is optimal
under all market regimes. We formalize this via a surrogate objective for the corrected
center $\hat{y}^*_t$:
%
\begin{equation}
\hat{y}^*_t = \arg\min_{c}
\left\{ \lambda_t\,(c - \hat{y}_t)^2 + (1-\lambda_t)\,(c - a_t)^2 \right\}
\end{equation}
%
whose closed-form solution is the convex combination:
%
\begin{equation}
\hat{y}^*_t = \lambda_t\,\hat{y}_t + (1-\lambda_t)\,a_t
\end{equation}
%
GSPQR sets $\lambda_t = G(r_t, \delta_t)$, where $\delta_t = \hat{y}_t - a_t$ is the
disagreement between predictors. The choice of $\lambda_t$ is thus data-driven and
locally adaptive, rather than fixed across all samples.

\subsection{Geometric Soft-Pullback Mechanism}
\label{sec:pullback}

\paragraph{Spiral anchor.}
We represent the prediction state as a 2D anchor point $\mathbf{p}_t = (a_t, r_t)$,
mapping to polar coordinates $(\rho_t, \theta_t)$ where
$\rho_t = \sqrt{a_t^2 + r_t^2}$ and $\theta_t = \arctan(r_t/a_t)$.
This representation is inspired by the logarithmic spiral $\rho = Ae^{B\theta}$, in
which the radial coordinate encodes prediction uncertainty and angular displacement
encodes directional signal. The soft-pullback corresponds to a radial displacement
along this curve toward the mean predictor, hence the name.

\paragraph{Uncertainty gate.}
We define the gating function as:
%
\begin{equation}
G(r_t,\, \delta_t) = \exp\!\left(-\beta \cdot \frac{|\delta_t|}{r_t}\right)
\end{equation}
%
The ratio $\eta_t = |\delta_t|/r_t$ serves as a \textit{signal-to-noise ratio} (SNR):
it measures the normalized disagreement between the mean and median estimates relative
to local prediction uncertainty. The gate exhibits three interpretable regimes:

\begin{itemize}
    \item \textbf{Inverse sensitivity.} When $r_t$ is large, $\eta_t \to 0$ and
    $G \to 1$, so $\hat{y}^*_t \to \hat{y}_t$. High uncertainty causes the mechanism
    to fall back on the more robust MSE estimate.
    \item \textbf{Adaptive recalibration.} When $r_t \to 0$, $G \to 0$ and
    $\hat{y}^*_t \to a_t$. High model confidence suppresses the MSE correction,
    preserving the directional signal from the median predictor.
    \item \textbf{Spiral logic.} The exponential form mirrors the geometry of the
    logarithmic spiral: correction strength transitions smoothly and nonlinearly with
    prediction confidence, rather than switching abruptly.
\end{itemize}

\begin{figure}[htbp]
    \centering
    \includegraphics[width=0.48\textwidth]{outputs/paper_plots/figure5_heatmap.png}
    \hfill
    \includegraphics[width=0.48\textwidth]{outputs/paper_plots/figure6_gate.png}
    \caption{GSPQR Decision Space Heatmap at $\beta=1.0$ (left) and Gating Function Continuous Behavior (right).}
    \label{fig:mechanism}
\end{figure}

\paragraph{Corrected interval.}
The final prediction interval is:
%
\begin{equation}
\hat{C}_t^* = \bigl[\hat{y}^*_t - r_t,\;\hat{y}^*_t + r_t\bigr]
\end{equation}
%
Interval width $2r_t$ is preserved by construction --- all improvements in Winkler
Score arise from better center placement, not artificial narrowing.

\paragraph{Special cases.}
When $\beta = 0$, $G \equiv 1$ and $\hat{y}^*_t = \hat{y}_t$, recovering the MSE
baseline (Snail-0). As $\beta \to \infty$, $G \to 0$ and $\hat{y}^*_t \to a_t$,
recovering the Q50-only baseline (Snail-$\infty$).

\subsection{Composite Scoring and $\beta$ Selection}

We select $\beta$ on a held-out validation set (2022 Q2--Q4) using a composite score
that jointly evaluates interval quality and calibration:
%
\begin{equation}
\text{Score} = \bar{W} + 10 \cdot \max(0,\, \text{CE} - 0.05)
\end{equation}
%
where $\bar{W}$ is the mean Winkler Score. The penalty activates only when coverage
deviation exceeds 5 percentage points, encouraging calibration while minimizing width.
We search $\beta \in \{0, 0.5, 1.0, 2.0, 5.0\}$ and select $\beta^* = 1.0$.

%%%%%%%%%%%%%%%%%%%%%%%%%%%%%%%%%%%%%%%%%%%%%%%%%%%
\section{Experiments}
%%%%%%%%%%%%%%%%%%%%%%%%%%%%%%%%%%%%%%%%%%%%%%%%%%%

\subsection{Data and Features}

We evaluate on HS300 constituent stocks with the dataset partitioned as in
Table~\ref{tab:splits}. The validation window (2022 Q2--Q4) is deliberately chosen
to coincide with a period of pronounced regime shift in the A-share market, providing
a stress-tested environment for $\beta$ selection.

\begin{table}[h]
\centering
\caption{Dataset splits.}
\label{tab:splits}
\begin{tabular}{lll}
\toprule
Split & Period & Purpose \\
\midrule
Training    & 2019-01-01 -- 2021-12-31 & Model parameter learning \\
Calibration & 2022-01-01 -- 2022-03-31 & CP baseline calibration \\
Validation  & 2022-04-01 -- 2022-12-31 & $\beta$ selection \\
Test        & 2023-01-01 -- 2024-12-31 & Out-of-sample evaluation \\
\bottomrule
\end{tabular}
\end{table}

We construct 12 features spanning four categories: momentum (\texttt{mom\_20d},
\texttt{mom\_60d}, \texttt{mom\_12m\_minus\_1m}), volatility (\texttt{vol\_60d\_res}),
valuation (\texttt{sp\_ratio}), and liquidity (\texttt{turn\_20d}). Each of these six
features is additionally transformed into a cross-sectional rank, yielding 12 features
in total. All features are normalized via rolling Z-score with a 60-day window.

\subsection{Baselines and Evaluation Metrics}

We compare against four baselines at $\alpha = 0.2$:
\begin{itemize}
    \item \textbf{Residual (Base).} Symmetric intervals $\hat{y}_t \pm 1.28\sigma$
    using calibration-set residual standard deviation. Assumes Gaussian, stationary noise.
    \item \textbf{CP (EnbPI).} Conformal prediction \citep{xu2021} calibrated on
    2022 Q1.
    \item \textbf{Snail-0 (MSE-base).} GSPQR at $\beta = 0$; reduces to MSE-centered
    QR intervals.
    \item \textbf{Snail-$\infty$ (Q50-base).} GSPQR at $\beta \to \infty$; reduces to
    Q50-centered intervals.
\end{itemize}

We report Actual Coverage (AC), Coverage Error CE, Winkler Score (primary metric),
interval Width, MAE of the corrected center, and RankIC (Spearman rank correlation
with realized returns).

\subsection{Main Results}

Table~\ref{tab:main} reports test-set performance ($n = 120{,}512$).

\begin{table}[h]
\centering
\caption{Test-set evaluation on A-share HS300 (2023--2024). Target coverage: 80\%.
$\downarrow$ lower is better, $\uparrow$ higher is better.}
\label{tab:main}
\small
\begin{tabular}{lcccccc}
\toprule
Method & AC & CE $\downarrow$ & Winkler $\downarrow$ & Width & MAE $\downarrow$ & RankIC $\uparrow$ \\
\midrule
Residual (Base)           & 0.9396 & 0.1396 & 0.5982 & 0.5239 & 0.1016 & $-$0.0081 \\
CP (EnbPI)                & 0.8904 & 0.0904 & 0.5083 & 0.3851 & 0.0918 & 0.0230 \\
Snail-0 (MSE-base)        & 0.8407 & 0.0407 & 0.4794 & 0.3197 & 0.0930 & 0.0024 \\
Snail-0.5                 & 0.8418 & 0.0418 & 0.4787 & 0.3197 & 0.0926 & 0.0072 \\
\textbf{GSPQR ($\beta^*=1.0$)} & \textbf{0.8423} & \textbf{0.0423} & \textbf{0.4784} & \textbf{0.3197} & \textbf{0.0922} & 0.0117 \\
Snail-2.0                 & 0.8433 & 0.0433 & 0.4784 & 0.3197 & 0.0917 & 0.0188 \\
Snail-5.0                 & 0.8449 & 0.0449 & 0.4789 & 0.3197 & 0.0909 & 0.0306 \\
Snail-$\infty$ (Q50-base) & 0.8457 & 0.0457 & 0.4815 & 0.3197 & 0.0897 & 0.0372 \\
\bottomrule
\end{tabular}
\end{table}

Several observations merit attention. First, the Residual baseline fails substantially
(CE $= 0.1396$, Winkler $= 0.5982$), confirming that Gaussian assumptions are
inappropriate for A-share return distributions. Second, CP (EnbPI) achieves adequate
coverage but at the cost of significantly wider intervals (Width $= 0.3851$, a 20.5\%
excess over QR-based methods). Third, all Snail variants share the same width of 0.3197,
confirming that performance gains arise solely from improved center placement.

\begin{figure}[htbp]
    \centering
    \includegraphics[width=0.7\textwidth]{outputs/paper_plots/figure4_pareto.png}
    \caption{Interval Quality Pareto Frontier: Mean Interval Width vs. Coverage Error.}
    \label{fig:pareto_efficiency}
\end{figure}

The primary finding is that \textbf{GSPQR ($\beta^* = 1.0$) achieves the lowest Winkler
Score of 0.4784}, a 5.9\% improvement over EnbPI and a 0.86\% MAE reduction relative
to Snail-0 --- reaching nearly 1\% enhancement in point forecasting accuracy despite no
direct MSE optimization. The improvement over EnbPI is statistically significant at
$p \ll 0.001$ (paired t-test, $t = -68.5$, $n = 120{,}512$).

\begin{figure}[htbp]
    \centering
    \includegraphics[width=0.9\textwidth]{outputs/paper_plots/figure2_interval_comparison.png}
    \caption{Empirical Interval Comparison: GSPQR vs. EnbPI over a sampled time series.}
    \label{fig:interval_comparison}
\end{figure}

\subsection{$\beta$ Parameter Analysis}

The Winkler Score as a function of $\beta$ exhibits a \textbf{convex profile} with a
minimum at $\beta^* \in [1.0, 2.0]$. At $\beta = 0$, the MSE center has poor ranking
ability (RankIC $= 0.0024$). As $\beta$ increases, RankIC rises monotonically (to
0.0372 at $\beta = \infty$), but coverage error also increases, eventually degrading
the Winkler Score. 

\begin{figure}[htbp]
    \centering
    \includegraphics[width=0.8\textwidth]{outputs/paper_plots/figure3_beta_sensitivity.png}
    \caption{Sensitivity Analysis of $\beta$: The Trade-off Frontier between Winkler Score and RankIC.}
    \label{fig:beta_tradeoff}
\end{figure}

The sweet spot at $\beta^* = 1.0$ captures the Q50 directional
signal before the coverage penalty dominates. This convex profile confirms that GSPQR
navigates a well-defined Pareto frontier between interval calibration and
cross-sectional ranking ability.

\subsection{Statistical Significance}

We conduct a paired t-test over per-sample Winkler Scores. For each test observation,
we compute $d_i = W_i^{\text{GSPQR}} - W_i^{\text{EnbPI}}$, yielding:
%
\begin{equation}
t = \frac{\bar{d}}{s_d / \sqrt{n}} = -68.5 \qquad (p \ll 0.001,\ n = 120{,}512)
\end{equation}
%
The large negative $t$-statistic confirms that GSPQR produces systematically lower
Winkler Scores across the entire test distribution, not merely on average.

%%%%%%%%%%%%%%%%%%%%%%%%%%%%%%%%%%%%%%%%%%%%%%%%%%%
\section{Discussion}
%%%%%%%%%%%%%%%%%%%%%%%%%%%%%%%%%%%%%%%%%%%%%%%%%%%

\subsection{Why Does the Soft-Pullback Work?}

The effectiveness of GSPQR stems from two complementary properties of A-share return
distributions. First, the conditional median and mean carry non-redundant information:
the median is robust to heavy tails while the mean encodes expected value across all
outcomes. The soft-pullback combines them in a locally adaptive ratio. Second, the
uncertainty gate provides automatic regime adaptation: during high-volatility periods
$r_t$ expands, attenuating the correction and falling back on the stable MSE estimate;
during trend-following regimes $r_t$ contracts, allowing the Q50 anchor to capture
cross-sectional directional signal.

\begin{figure}[htbp]
    \centering
    \includegraphics[width=0.9\textwidth]{outputs/paper_plots/figure1_regime_adaptation.png}
    \caption{Adaptive Correction Path Across Different Market Regimes.}
    \label{fig:regime_adaptation}
\end{figure}

\subsection{The $\beta$ Trade-off as a Convex Problem}

The convex Winkler profile reveals a fundamental tension: calibration and ranking are
partially conflicting objectives. The $\beta$ parameter traces a Pareto frontier between
them, and the composite scoring criterion selects the optimal balance. This framing
positions GSPQR not as a heuristic but as a principled solution to a well-defined
optimization problem over the space of interval center estimators.

\subsection{Bayesian Analogy}

The gate function $G(r_t, \delta_t)$ admits a Bayesian interpretation. In a Bayesian
Model Averaging framework, model weights are proportional to likelihood. Here, the
``likelihood'' of the mean predictor is measured by the inverse SNR $|\delta_t|/r_t$.
GSPQR effectively performs adaptive shrinkage toward the median, robustifying estimation
precisely when the signal is most ambiguous --- a behavior consistent with empirical
Bayes shrinkage estimators.

\subsection{Limitations}

Several limitations should be acknowledged. First, evaluation is restricted to a
single universe (HS300) over a two-year test window; generalization to the CSI 500 or
small-cap stocks remains to be verified. Second, the parsimonious feature set (12
features) ensures reproducibility but may interact differently with richer alternative
data. Third, the symmetric interval assumption may be suboptimal under skewed
distributions; extending GSPQR to asymmetric intervals is a natural direction for
future work.

%%%%%%%%%%%%%%%%%%%%%%%%%%%%%%%%%%%%%%%%%%%%%%%%%%%
\section{Conclusion}
%%%%%%%%%%%%%%%%%%%%%%%%%%%%%%%%%%%%%%%%%%%%%%%%%%%

We have presented GSPQR, a geometric soft-pullback mechanism for quantile interval
prediction in equity markets. By formalizing interval re-centering as an adaptive
SNR-aware optimization and implementing it via an exponential uncertainty gate inspired
by logarithmic spiral geometry, GSPQR achieves prediction intervals that are
simultaneously well-calibrated and informationally efficient. On the HS300 test set
(2023--2024, $n = 120{,}512$), GSPQR achieves a Winkler Score of 0.4784, outperforming
EnbPI by 5.9\% ($t = -68.5$, $p \ll 0.001$) while reducing interval width by 17.0\%.

These results suggest that the standard separation between interval construction and
point prediction is suboptimal in cross-sectional equity settings. GSPQR provides a
lightweight, interpretable bridge between the two with no increase in model complexity.
Future work will explore asymmetric interval extensions, integration with deep sequence
models, and evaluation across broader A-share universes.


%%%%%%%%%%%%%%%%%%%%%%%%%%%%%%%%%%%%%%%%%%%%%%%%%%%
\section*{Acknowledgements}
%%%%%%%%%%%%%%%%%%%%%%%%%%%%%%%%%%%%%%%%%%%%%%%%%%%

The author acknowledges the assistance of multiple large language models in the development and refinement of this work.This work was supported by AI-assisted research tools, including ChatGPT, Claude, Gemini, and GLM, which contributed to idea exploration, code review, and writing refinement.

In particular, ChatGPT (OpenAI), Claude (Anthropic), Gemini (Google DeepMind), and GLM (Zhipu AI) were used as interactive research assistants for idea clarification, language polishing, structural feedback, and preliminary code/proof checking.

All core methodological design, experimental implementation, and final academic judgments remain the sole responsibility of the author.

\bibliographystyle{apalike}
\bibliography{references}

\end{document}
